\section{Finiteness in a fixed band}\label{sec:exact-band}

\begin{corollary}[Finiteness for the infimum]\label{cor:entered-type-attainment}
For each \(\delta>0\), only finitely many oriented nonstandard smooth
homotopy four-sphere types have \(I^+(M)\le 2T-\delta\).
\end{corollary}
\begin{proof}
Even at equality this infimum supplies an entering class, since
\(\delta>0\). Apply Theorem~\ref{thm:attainment} to take a positive
Bach-flat class on the same type at the infimum itself. Its total
Weyl integral is uniformly bounded. In our Schouten normalization,
\[
       \int_M\sigma_2(g^{-1}A)\,dV=(2T-e)/2\ge\delta/2 .
\]
These are the hypotheses of the finiteness theorem of Chang, Qing and
Yang \cite[Theorem B]{ChangQingYang2006Gap}. Their Schouten tensor in
the statement is \(2A\), so their lower bound is \(2\delta\).
Each unoriented smooth type has at most two oriented types.
\end{proof}

Consequently, for any sequence of pairwise distinct nonstandard entered types, their infima must tend to \(2T\) from below. On the positive
Bach-flat classes in a fixed nonstandard band \(e\le2T-\delta\), the spectral deficit
\(36\int L^2\,dV-Y^2\) has a strictly positive attained minimum whenever
the family is nonempty. The least exact value on each type is isolated
among exact values. These refinements use the compactness or strict
energy drop of the marked-neck theorem
\cite[First fixed-type critical value]{Baker2026MarkedNecks}
\cite[The spectral deficit of exact critical classes]{Baker2026FixedDescent}.
Neither finiteness nor isolation excludes an individual nonstandard type.
